\honest{Taboo Your Words}{Eliezer Yudkowsky}{2008}

In the game Taboo you must get your partner to guess a word without saying it or five related words. Given ``baseball,'' I at once think of ``An artificial group conflict in which you use a long wooden cylinder to whack a thrown spheroid.'' I have practiced blanking words out of my mind for years.

Yesterday, I say, replacing ``human'' by its definition showed that the classic syllogism about Socrates is empty. The definition I used already contained ``mortal.'' I do not consider a definition that leaves it out. \nb{John Stuart Mill made the same charge against the syllogism in 1843.}

Then the falling tree again: one person says it makes a sound, the other says it does not. Replace ``sound'' with what each of them means, vibrations in the air or a heard experience, and the conflict disappears. If ``vibrations'' comes into dispute, taboo that too, and then ``wave,'' and so on.

The reverse can also happen: two people use the same word and expect different things. Banning the word would show the difference. This is the case where the method would find a real disagreement, and I give no real example of it.

Instead I give predictions. Futurists who agree that we will have artificial intelligence in thirty years, I say, ``might find quite a conflict of expectations'' once the term is banned. I point to a collection of 71 definitions of intelligence by Shane Legg and Marcus Hutter. \nb{They write that there is ``still no standard definition of intelligence,'' and also that many of the definitions ``share many common features.''} Religious people who seem united: ban ``God'' and ``faith,'' and ``mostly they won't be able to answer at all.'' Nobody has been asked.

Near the end comes the thesis. In philosophy, the first line of defense ``is not to define your problematic terms, but to see whether you can think without using those terms at all.'' I add that you must also ban ``any of their short synonyms,'' and not ``invent a new word to use instead.'' Instead: ``Describe outward observables and interior mechanisms.''

On free will, I say, most philosophers would advise defining the term; I advise describing what people do or do not have without it. ``If you want to try this at home,'' I add, also avoid ``choose,'' ``act,'' ``decide,'' ``determined'' and ``responsible.'' I do not do the exercise myself. A disagreement over whether people are morally responsible for what they do is where a real dispute about free will would show up. I ban ``responsible'' too, and I do not show how that dispute would be stated without it.

I conclude: ``This is one of the nonstandard tools in my toolbox, and in my humble opinion, it works way way better than the standard one.'' I have not shown it working on any open question.
